%% chapter3.tex: System Design and Implementation

\chapter{System Design and Implementation}
\label{chap:implementation}

This chapter describes the design and implementation of the routing system developed in this
dissertation. The system is built on top of the TDMA framework of
Almeida~\cite{almeida2023synchronization} and retains the RA-TDMAs+ protocol stack
described in Chapter~\ref{chap:background}. The central departures from the previous work of
Morais~\cite{morais2024synchronization} are the adoption of a Layer~3 forwarding model, the
separation of the communication stack into a UDP control plane and a TCP data plane, and the
addition of per-link quality monitoring to the routing decision. The following sections describe
each design decision and its implementation in detail.


%% -----------------------------------------------------------------------
\section{Protocol Architecture: Separating Control and Data Planes}
\label{sec:proto-split}

The original TDMA framework used UDP datagrams for all inter-node communication:
both the MATRIX topology beacons and the application data packets were transmitted through
the same UDP socket. UDP is appropriate for the MATRIX control traffic because the periodic
nature of the beacon means that a lost packet is simply replaced by the next one, one round
later, without any lasting effect on the shared topology state. It is, however, unsuitable
as the sole transport for video streaming: each lost datagram causes visible artefacts at
the destination, and UDP provides no retransmission mechanism to recover lost frames.

Prior work by Morais~\cite{morais2024synchronization} tested TCP at the application layer
as an alternative video transport. While TCP eliminates packet-loss artefacts through its
retransmission mechanism, that work observed that default TCP socket settings introduce
unacceptable delays for real-time streaming over TDMA: Nagle's algorithm buffers
outgoing data until a full TCP segment accumulates or an in-flight segment is acknowledged,
introducing latency of up to 200\,ms per hop, and default keep-alive timers of two hours
leave dead connections undetected for the duration of a mission. The present work moves TCP
connection management into the TDMA framework itself, with per-peer sockets explicitly
configured to avoid these issues, as described in the following subsection.

To address this, the present work separates the communication stack into two planes, each
with a transport protocol matched to its requirements.

The \textbf{control plane} continues to use UDP. MATRIX topology beacons are multicast
datagrams broadcast at the start of each slot. Their content is redundant across rounds (the
same connectivity information is retransmitted every round), so occasional packet loss has no
lasting consequence. The stateless, broadcast nature of the beacon is also a better fit for UDP
than for a connection-oriented protocol.

The \textbf{data plane} uses TCP. Application data packets, encapsulated as \texttt{MSG\_DATA}
frames, are delivered over dedicated per-peer TCP connections. TCP provides guaranteed
in-order delivery through its retransmission and acknowledgement mechanism, ensuring that
video frames arrive intact at the destination regardless of transient link quality changes.
Each node maintains an array of TCP socket file descriptors indexed by peer identifier,
\texttt{tcp\_sockfd[MAX\_NODES+1]}, so that each pair of communicating nodes has its own
dedicated stream with no interference from control traffic.

\subsection{TCP Framing}

Because TCP is a byte-stream protocol with no inherent message boundaries, each logical
message must carry its own length information. Every outgoing \texttt{MSG\_DATA} frame is
prefixed with a four-byte network-order integer that encodes the total length of the following
content. Figure~\ref{fig:frame-structure} shows the structure of both message types used in
the system.

\begin{figure}[H]
\centering
\includegraphics[width=0.90\textwidth]{frame_structure}
\caption{Frame structure for MATRIX Beacon (UDP, Type~1) and MSG\_DATA (TCP, Type~2).}
\label{fig:frame-structure}
\end{figure}

The receiver calls \texttt{recv()} with the \texttt{MSG\_WAITALL} flag twice in sequence: first
to read exactly four bytes and determine the message length, then to read exactly that many
bytes to obtain the complete packet. This two-step receive guarantees that the application
always processes whole, aligned messages, regardless of how TCP segments the byte stream
at the network level.

\subsection{TCP Socket Configuration}
\label{subsec:tcp-socket-config}

The TCP sockets are configured with a set of options that are essential for correct and
low-latency operation. Without each option the system would either introduce unacceptable
delays, fail to detect broken connections, or block indefinitely and disrupt the TDMA slot
schedule. Code Block~\ref{lst:tcp-opts} shows the options applied to every peer connection.

\begin{lstlisting}[language=C,
    caption={TCP socket options applied to each peer connection},
    label={lst:tcp-opts}]
setsockopt(fd, IPPROTO_TCP, TCP_NODELAY,   &one,      sizeof(one));
setsockopt(fd, SOL_SOCKET,  SO_KEEPALIVE,  &one,      sizeof(one));
setsockopt(fd, IPPROTO_TCP, TCP_KEEPIDLE,  &(int){5}, sizeof(int));
setsockopt(fd, IPPROTO_TCP, TCP_KEEPINTVL, &(int){2}, sizeof(int));
setsockopt(fd, IPPROTO_TCP, TCP_KEEPCNT,   &(int){3}, sizeof(int));
setsockopt(fd, SOL_SOCKET,  SO_SNDTIMEO,   &tv_snd,   sizeof(tv_snd));
setsockopt(fd, SOL_SOCKET,  SO_RCVTIMEO,   &tv_rcv,   sizeof(tv_rcv));
\end{lstlisting}

\textbf{\texttt{TCP\_NODELAY}} disables Nagle's algorithm, which accumulates outgoing data
in a buffer until a full TCP segment is ready or a previous segment has been acknowledged.
Without this option, small video packets would be held by the kernel for up to 200\,ms
waiting to be bundled with subsequent data. Since the transmitter thread must complete all
sends within its TDMA slot window (50\,ms for a three-node network), a packet delayed by
Nagle's buffering could miss the slot entirely and wait a full round before retransmission,
introducing one frame period of avoidable latency per hop.

\textbf{\texttt{SO\_KEEPALIVE}} enables TCP-level keep-alive probes at the kernel level.
Without it, if a peer node moves out of radio range or crashes without sending a TCP FIN or
RST, the local socket would remain in the \texttt{ESTABLISHED} state indefinitely. The routing
module would continue directing packets toward the dead peer, those packets would queue in
the transmission buffer, and no failure would ever be signalled to the application. The network
would appear to the operator as if packets were being forwarded normally, while in reality they
were being silently discarded.

\textbf{\texttt{TCP\_KEEPIDLE}} sets the idle time before the first keep-alive probe is sent.
The Linux default is 7200 seconds (two hours). In a robotic scenario where a node can move
out of range in a matter of seconds, a two-hour detection timeout is entirely unusable. Setting
this to 5 seconds ensures that an unresponsive peer is probed within five seconds of the last
received data.

\textbf{\texttt{TCP\_KEEPINTVL}} and \textbf{\texttt{TCP\_KEEPCNT}} control how often probes
are repeated and how many failures are tolerated before the connection is declared dead. With
the values used (interval of 2\,s, count of 3), the worst-case detection time after the idle
period expires is $5 + 3 \times 2 = 11$\,s. Under the Linux defaults (interval of 75\,s, count
of 9) this would be $7200 + 9 \times 75 = 7875$\,s, more than two hours, which is
impractical for a mobile network.

\textbf{\texttt{SO\_SNDTIMEO}} applies a one-second timeout to send operations. Without it,
if the TCP send buffer fills because the peer is slow or unreachable but not yet detected as
dead, the \texttt{send()} call inside the transmitter thread would block indefinitely. Because the
transmitter thread must release the channel before the TDMA slot boundary, an indefinitely
blocking send would cause it to overstay the slot and transmit into the allocation of the following node, corrupting slot boundaries for all nodes in the network.

\textbf{\texttt{SO\_RCVTIMEO}} applies a two-second timeout to receive operations. Without
it, the per-peer receive thread would block on \texttt{recv()} indefinitely if the peer stops
sending data without closing the connection cleanly, a condition known as a half-open
connection. The thread would remain stuck, the PDR module would stop receiving sequence
number updates for that peer, and stale link quality data would continue to influence routing
decisions long after the link had effectively failed.


%% -----------------------------------------------------------------------
\section{Architecture Overview}
\label{sec:arch-overview}

Figure~\ref{fig:framework-miguel} presents the complete architecture of the system developed
in this work. The TDMA Core and its four internal threads (TX, RX, Event Handler, and
Config), the Topology-Tracking module, and the RA-TDMAs+ synchronisation module are
retained without modification from the framework of Almeida~\cite{almeida2023synchronization}.
The Routing module and the TCP Manager are the central new components added to the
framework.

\bigskip
\begin{figure}[H]
\centering
\includegraphics[width=0.98\textwidth]{frameworkmiguel}
\caption{Final architecture of the TDMA framework extended with the Routing module and TCP Manager.}
\label{fig:framework-miguel}
\end{figure}
\bigskip

The Topology-Tracking module continues to maintain the connectivity matrix and runs
Prim's algorithm to compute the Minimum Spanning Tree, storing the result in a
shared-memory region. It notifies the TDMA Core Event Handler via a local socket whenever
the slot assignment changes due to a node joining or leaving, carrying the updated slot
identifier so the Core can reorganise the round before the next transmission.

The Routing module reads two inputs from shared memory: the Spanning Tree produced by
the Topology-Tracking module, and the per-link quality scores maintained by the PDR
sub-module. Because node identifiers in the framework correspond directly to the last octet
of the IP address of each node, no separate IP address distribution is required; the Routing
module derives every destination IP from the node identifiers already present in the Spanning
Tree. This contrasts with the Layer~2 approach of Morais~\cite{morais2024synchronization},
which required distributing MAC addresses through the MATRIX beacon because MAC
addresses cannot be derived from node identifiers. After computing next-hop decisions
through a Depth-First Search over the Spanning Tree, the Routing module installs the
resulting routes into the Linux kernel via Netlink, writing entries into two routing tables:
the main table and a secondary policy table (table~200). The two-table approach allows
application traffic and TDMA control traffic to follow separate forwarding paths, preventing
routing loops between the wireless interface and the TUN interface.

The TCP Manager groups three threads responsible for the data-plane connections to peer
nodes: the TCP acceptor thread, which accepts incoming connections from neighbours; the
TCP keepalive thread, which monitors active connections and reconnects dead peers; and one
TCP receiver thread per connected peer, which processes incoming \texttt{MSG\_DATA} frames.
All three threads interact directly with the Routing module through function calls within the
same process: the TX thread queries the Routing module for the correct TCP socket before
each transmission, and the per-peer receiver threads consult it to decide whether an arriving
packet should be delivered locally via \texttt{tun0} or forwarded to the next hop.


%% -----------------------------------------------------------------------
\section{Layer 3 Routing Design}
\label{sec:layer3-design}

The forwarding mechanism operates at the network layer (Layer~3), in contrast to the
Layer~2 ARP-manipulation approach of Morais~\cite{morais2024synchronization}.
Figure~\ref{fig:l2-vs-l3} illustrates the key structural difference between the two
approaches.

\begin{figure}[H]
\centering
\includegraphics[width=0.95\linewidth]{layer2v3}
\caption{Layer~2 vs Layer~3 routing: relay nodes process up to Data Link in L2 and up
to Network in L3.}
\label{fig:l2-vs-l3}
\end{figure}

\subsection{Layer~2 Approach (Morais, 2024)}

The work of Morais~\cite{morais2024synchronization} successfully demonstrated multi-hop
routing within the RA-TDMAs+ framework by operating at the Data Link layer. Relay nodes
forward packets by rewriting the destination MAC address of each frame to that of the next
hop and retransmitting on the wireless interface, leaving the IP header untouched throughout
the entire path. This approach proved effective in the experimental context of that work,
requiring no changes to the kernel IP stack and keeping the implementation compact.

To enable MAC-based forwarding, Morais extended the MATRIX beacon to distribute the MAC address of each node across the mesh, a clean solution that kept all necessary information
within the existing beacon infrastructure. A notable advantage of this design is that
topology changes require only a lightweight update to the ARP table entries, with no
interaction with the kernel routing subsystem. This makes route recomputation fast and
with low overhead, which is particularly beneficial in dynamic mesh topologies where nodes
join and leave frequently.

\subsection{Layer~3 Approach (This Work)}

The present work adopts a different design strategy, operating at the Network layer through
a TUN virtual interface (\texttt{tun\textit{id}}). The primary motivation was the adoption
of TCP as the data transport: because TCP connections are identified by IP addresses and
port numbers, the forwarding path must be visible at the IP layer. Combining TCP with
Layer~2 MAC rewriting would require each relay node to maintain per-flow TCP state, which
would be architecturally complex.

By operating at Layer~3, the system gains compatibility with the Linux kernel networking
stack: \texttt{iptables}, policy routing tables, and \texttt{traceroute} all function
correctly over the relay path. This is what enables the two-table traffic separation
described in Section~\ref{sec:kernel-config} and the use of Netlink for route management.
The trade-off is that each topology change requires installing and removing host routes via
Netlink, which is more expensive than an ARP table update. This cost is mitigated by the
use of Netlink sockets directly (approximately 50\,\textmu s per route) rather than
spawning \texttt{ip route} as a child process, but it remains a higher overhead than the
Layer~2 approach for frequent topology changes.


%% -----------------------------------------------------------------------
\section{Thread Architecture}
\label{sec:threads}

The system runs seven fixed threads per node plus one dynamic TCP receiver thread per
connected peer. Table~\ref{tab:threads} summarises their roles. The Receiver, Transmitter,
and Event Handler threads are retained from the TDMA framework of
Almeida~\cite{almeida2023synchronization}; the remaining threads are new contributions of
this work.

\begin{table}[H]
\centering
\caption{Threads running per node and their responsibilities.}
\label{tab:threads}
\begin{tabular}{lll}
\toprule
\textbf{Thread} & \textbf{Entry point} & \textbf{Role} \\
\midrule
Receiver          & \texttt{receiver\_loop()}        & Receives MATRIX beacons via UDP \\
Transmitter       & \texttt{tx\_loop()}               & Sends MATRIX (UDP) and MSG\_DATA (TCP) during slot \\
TUN reader        & \texttt{tun\_reader\_loop()}      & Reads IP packets from \texttt{tun0}, pushes to tx\_queue \\
Event handler     & \texttt{event\_handler\_loop()}   & Processes topology events, triggers re-routing \\
TCP acceptor      & \texttt{tcp\_accept\_loop()}      & Accepts incoming TCP connections from peers \\
TCP keepalive     & \texttt{tcp\_keepalive\_loop()}   & Monitors peer connections, reconnects dead ones \\
TCP RX (per peer) & \texttt{tcp\_rx\_peer\_loop()}    & Receives MSG\_DATA from one connected peer \\
\bottomrule
\end{tabular}
\end{table}

The transmitter thread is the timing-critical component. It wakes at the beginning of the
assigned TDMA slot of the node, drains the transmission queue, and sends each queued packet
over the appropriate TCP connection to the next-hop peer determined by the routing module.
The TUN reader thread runs continuously, capturing outgoing IP packets from \texttt{tun0} and
pushing them into the transmission queue with the destination peer ID, decoupling packet
capture from the slot-timed transmission.

The TCP acceptor thread executes a blocking \texttt{accept()} call and, upon each new
connection, records the peer socket in the \texttt{tcp\_sockfd[]} array and spawns a dedicated
\texttt{tcp\_rx\_peer\_loop} thread for that peer. Each per-peer receive thread calls
\texttt{recv()} with \texttt{MSG\_WAITALL} in a tight loop, processing one complete framed
message per iteration: it first reads the four-byte length prefix, then reads exactly that many
bytes for the payload. On delivery, it passes the packet to the routing layer, which either
injects it into \texttt{tun0} for local consumption or forwards it toward the next hop.

The TCP keepalive thread complements the kernel-level keep-alive mechanism with an
application-level health check. Every second it probes each peer socket with a non-blocking
\texttt{recv()} call using \texttt{MSG\_PEEK | MSG\_DONTWAIT}. A return value of zero
indicates a graceful connection close; a negative return with a non-recoverable error indicates
a reset or broken link. In either case the socket is closed, the \texttt{tcp\_sockfd[]} entry is
set to $-1$, and a reconnection attempt is scheduled. A race condition between the keepalive
thread closing a socket and a receive thread still using it is avoided by copying the file
descriptor under the TCP mutex before use and verifying that the descriptor has not been
replaced before closing it.

\subsection{Synchronisation Primitives}

Four synchronisation primitives protect shared state accessed by multiple threads:

\begin{itemize}
  \item \texttt{tcp\_mutex}: protects \texttt{tcp\_sockfd[]} against concurrent reads by the
    keepalive thread and writes by the acceptor thread.
  \item \texttt{tx\_queue->lock}: an internal mutex protecting the circular buffer head and
    tail indices during push and pop operations.
  \item \texttt{g\_matrix\_mutex}: protects the global connectivity matrix during updates by
    the receiver thread and reads by the routing and synchroniser modules.
  \item Event queue lock: a producer-consumer mutex protecting the event queue shared
    between the receiver thread (producer) and the event handler thread (consumer).
\end{itemize}


%% -----------------------------------------------------------------------
\section{Data Structures}
\label{sec:data-structures}

\subsection{Transmission Queue}

The transmission queue (\texttt{tx\_queue.c}) is a circular buffer that stores packets captured
from \texttt{tun0} until the transmitter thread can send them during the TDMA slot of the node.
Its parameters are a capacity of 64 packets (\texttt{TX\_QUEUE\_CAPACITY}) and a maximum
entry size of 2048 bytes (\texttt{TX\_QUEUE\_MAX\_PKT\_SIZE}). When the buffer is full, the
oldest packet at the head is discarded to make room for the new arrival (tail-drop policy),
and a warning is logged. Push and pop operations are non-blocking and mutex-protected.

The tail-drop policy is appropriate for video streaming: a packet that has been waiting in
the queue through one or more missed slots is less valuable to the receiver than a fresh
frame, so discarding old entries on overflow is preferable to discarding new ones.

\subsection{Routing Linked List}

The routing module retains the linked-list structure introduced by
Morais~\cite{morais2024synchronization}: primary list entries represent directly connected
neighbours, and each primary entry carries a secondary list of all nodes reachable through
that neighbour, as described in Section~\ref{sec:related-work}. The \texttt{routingList\_t}
type is unchanged. In the Layer~3 design no MAC address distribution is needed: the routing module derives each
destination IP address directly from its node identifier as the last octet of the
10.0.0.0/24 prefix, so the lookup result is a TCP socket index rather than a link-layer
address.

\subsection{Per-peer Quality State}

Each peer is tracked through two fields: its TCP socket file descriptor in
\texttt{tcp\_sockfd[]} and its link quality score maintained by \texttt{pdr.c}. The link
quality score is an integer in the range $[0, 100]$, updated after every ten received packets
and used as the edge weight in the Minimum Spanning Tree computation.


%% -----------------------------------------------------------------------
\section{Routing Algorithm}
\label{sec:routing-algo}

\begin{figure}[H]
\centering
\includegraphics[width=0.88\textwidth]{routingModule}
\caption{Execution cycle of the Routing module.}
\label{fig:routing-module}
\end{figure}

The routing algorithm follows the MST-based proactive approach identified in
Section~\ref{sec:routing-applicability}. The Routing module executes as a dedicated thread
driven by topology-change events. Figure~\ref{fig:routing-module} shows its execution cycle,
which repeats each time the mesh connectivity changes.

The cycle repeats each time the mesh connectivity changes and proceeds through a
fixed sequence of steps. The module first waits for an event, blocking on the
event queue until the Event Handler delivers a topology-change notification,
which is issued whenever the Topology-Tracking module detects a node joining or
leaving the mesh and the MATRIX connectivity matrix changes. On wake-up, it reads
the current set of active node identifiers and the Minimum Spanning Tree from the
shared-memory regions maintained by the Topology-Tracking module, both protected
by \texttt{g\_matrix\_mutex} to prevent a torn read during a concurrent update.
It then computes the next hop with a Depth-First Search (DFS) on the Spanning Tree
rooted at the local node: for each reachable destination the DFS records the
first-hop neighbour through which it is reached, and the set of destinations
associated with each first-hop neighbour is sorted by node identifier using
\texttt{qsort} to produce the secondary linked list for that primary entry. The
routing table is then updated by comparing the sorted DFS output against the
current linked-list routing table in a single linear scan, inserting nodes present
in the new output but absent from the list, removing nodes present in the list but
absent from the output, and leaving unchanged entries untouched, so that the number
of list operations is proportional to the size of the topology change rather than
the total network size. For every modified entry the module then updates the kernel
routes, installing or removing two host routes per destination, one in the main
routing table for application-originated traffic and one in policy table~200 for
relay-forwarded traffic, as described in detail in
Section~\ref{sec:kernel-config}. Finally, the module sleeps for a short settling
interval before processing the next event, which prevents route flapping during
rapid topology transients.

When a packet must be forwarded, \texttt{routing\_manager\_lookup()} traverses the primary
linked list and searches each secondary list for the destination identifier. The first matching
primary entry yields the next-hop peer ID, which indexes into \texttt{tcp\_sockfd[]} to obtain
the TCP socket for transmission.


%% -----------------------------------------------------------------------
\section{Kernel Network Configuration}
\label{sec:kernel-config}

Correct operation of the Layer~3 forwarding path requires a set of kernel network rules
established once at startup by the TUN module. This section describes the four configuration
groups and their interaction with the routing tables maintained by the Routing module.
Figure~\ref{fig:trafico} shows the two traffic paths and the routing tables that govern them.

\begin{figure}[H]
\centering
\includegraphics[width=0.95\linewidth]{trafico}
\caption{Separation of locally originated application traffic (via Table~100) and relayed transit traffic (via Table~200).}
\label{fig:trafico}
\end{figure}

Three routing tables are involved, each with a distinct role, and it is worth
stating them together before the detailed rules to avoid confusion. Table~100
holds a single default route through \texttt{tun\textit{id}} and exists only to
divert locally generated application traffic into the TUN interface, so that the
framework can capture and transport it. The main table holds the per-destination
application next-hop routes installed by the Routing module, which determine the
next hop toward each mesh destination for traffic originated by the local node.
Table~200 holds the per-destination host routes used when a relay node forwards
transit traffic back out on \texttt{wlan0}. The first is established once at
startup by the TUN module, while the main-table and table~200 entries are
maintained by the Routing module and updated on every topology change.

\subsection{IP Forwarding and Interface Setup}

Two kernel parameters are set at startup. \texttt{net.ipv4.ip\_forward} is enabled so
that the kernel forwards IP packets between interfaces; without it, packets arriving on
\texttt{tun\textit{id}} would be silently dropped instead of being routed towards
\texttt{wlan0}. \texttt{net.ipv4.conf.*.rp\_filter} is disabled on all relevant interfaces
because relay traffic injected through \texttt{tun\textit{id}} originates from the
10.0.0.0/24 prefix but exits on \texttt{wlan0}; the reverse-path check of the kernel would
reject these packets if left enabled.

The wireless interface \texttt{wlan0} is placed in IEEE~802.11 ad-hoc mode with a fixed
ESSID and channel, then assigned a physical address in the \texttt{172.20.10.0/28} subnet.
The TUN interface \texttt{tun\textit{id}} is assigned the logical address
\texttt{10.0.0.\textit{id}/24}, and a subnet route for \texttt{10.0.0.0/24} is added
pointing to it. The two address spaces are intentionally kept separate: the
\texttt{172.20.10.x} addresses are used exclusively for physical next-hop forwarding in
table~200, while the \texttt{10.0.0.x} namespace is the one exposed to applications and
to the routing algorithm.

\subsection{Application Traffic Redirection (Table~100)}

Application traffic addressed to the mesh (10.0.0.0/24) must be intercepted before it
reaches \texttt{wlan0} and redirected through \texttt{tun\textit{id}} so the TUN reader
thread can capture it. This is achieved through three iptables and policy routing rules.
First, an \texttt{iptables} mangle rule in the OUTPUT chain marks all outgoing
\texttt{wlan0} packets destined for \texttt{10.0.0.0/24} with \texttt{fwmark~1}. A second
rule immediately clears the mark for TDMA beacon packets, identified by their fixed UDP
source port (\texttt{7000\,+\,\textit{id}}), so that the control traffic of the framework itself
never enters the TUN path. Finally, a policy routing rule directs all packets with
\texttt{fwmark~1} to table~100, which holds a single default route via
\texttt{tun\textit{id}}.

Every marked packet exits via \texttt{tun\textit{id}} and appears at the corresponding
file descriptor. The TUN reader thread reads it, extracts the destination node identifier
from the IP destination address, and pushes the packet into the transmission queue. During
its TDMA slot, the transmitter thread pops the packet, queries the routing table for the
next-hop peer ID, and sends the framed packet over the corresponding TCP connection.

\subsection{Relay Forwarding Path (Table~200)}

When a relay node receives a TCP packet destined for a further node, it injects the IP
payload into \texttt{tun\textit{id}} via \texttt{tun\_write()}. From the perspective of the kernel
the packet appears as traffic arriving on the TUN interface and requiring IP
forwarding. A policy rule with priority~100, lower than the default table priority of
32766, directs all packets ingressing on \texttt{tun\textit{id}} to table~200, where
the Routing module installs a \texttt{/32} host route per destination. Each host route
specifies the physical \texttt{172.20.10.\textit{x}} address of the next-hop node as the
gateway and \texttt{wlan0} as the outgoing interface, so forwarded packets exit directly
onto the ad-hoc link without re-entering \texttt{tun\textit{id}}. Two \texttt{iptables}
FORWARD rules explicitly permit traffic in both directions between \texttt{tun\textit{id}}
and \texttt{wlan0}, since the default FORWARD policy of the Linux kernel is DROP and relay
traffic would otherwise be silently discarded.

\subsection{Kernel Route Installation via Netlink}

Route entries in both the main table and table~200 are installed and removed using Netlink
sockets (\texttt{AF\_NETLINK}, \texttt{NETLINK\_ROUTE}) rather than by spawning
\texttt{ip route} as a child process. Each Netlink message contains an \texttt{nlmsghdr}
header (type \texttt{RTM\_NEWROUTE} or \texttt{RTM\_DELROUTE}), an \texttt{rtmsg}
routing header specifying the address family, prefix length, table identifier, protocol,
scope, and route type, and a set of variable-length \texttt{rtattr} attributes carrying the
destination address (\texttt{RTA\_DST}), gateway (\texttt{RTA\_GATEWAY}), and output
interface index (\texttt{RTA\_OIF}). The full round-trip including the kernel ACK costs
approximately 50\,\textmu s per route, compared to approximately 3\,ms for an equivalent
\texttt{system("ip route ...")} call, a 60-fold reduction significant when updating routes
for the entire mesh at topology-change time.

Code Block~\ref{lst:netlink-api} shows the four API calls used by the Routing module.

\begin{lstlisting}[language=C,
    caption={Netlink route management API used by the Routing module},
    label={lst:netlink-api}]
/* topology change: install two routes per destination */
ip_route_add("10.0.0.3", "10.0.0.2", "tun1");         /* main table  */
ip_route_add_t("10.0.0.3", "172.20.10.3", "wlan0", 200); /* table 200 */

/* node departure: remove routes */
ip_route_del("10.0.0.3", "tun1");   /* main table entry */
ip_route_flush_table(200);           /* purge all table-200 entries */
\end{lstlisting}

\texttt{ip\_route\_add()} uses flags \texttt{NLM\_F\_CREATE | NLM\_F\_REPLACE}, so a
repeated call for the same destination silently replaces the existing entry without error.
At teardown, \texttt{ip\_route\_flush\_table(200)} issues a single
\texttt{system("ip route flush table 200")} call to purge all relay routes at once.


%% -----------------------------------------------------------------------
\section{Slot Data Control}
\label{sec:slot-control}

A limitation of the routing work of Morais~\cite{morais2024synchronization} was the
absence of any mechanism to limit the volume of data transmitted within a single TDMA slot.
Without such a mechanism, a burst of queued packets could cause the transmitter thread to
overstay its slot, causing its transmissions to overlap with the slot of the following node and
degrading synchronisation quality.

This work introduces a per-slot byte budget, \texttt{MAX\_BYTES\_SLOT}, computed from the
configured WiFi link rate and the slot duration. The transmitter thread tracks elapsed time
and stops dequeuing when the slot boundary is reached. Figure~\ref{fig:buffercontrol}
illustrates the mechanism: packets P1, P2 and P3 (blue) are sent during the slot; when
\texttt{slot\_end} is detected, packets P4 and P5 (red) are pushed back into the queue
and sent in the following round.

\begin{figure}[H]
\centering
\includegraphics[width=0.95\linewidth]{buffercontrol}
\caption{Slot data control: packets sent during the slot (blue) and pushed back on slot\_end (red).}
\label{fig:buffercontrol}
\end{figure}

\begin{lstlisting}[language=C,
    caption={Back-pressure mechanism in the transmitter thread (\texttt{tx\_loop})},
    label={lst:backpressure}]
while ((pkt = tx_queue_pop(node->tx_queue)) != NULL) {
    if (get_time_us() >= slot_end) {
        tx_queue_push(node->tx_queue, pkt->data,
                      pkt->len, pkt->dst_id);
        break; /* return packet to queue for next slot */
    }
    /* send packet over TCP to next-hop peer */
}
\end{lstlisting}

If the slot boundary is reached before the queue is empty, the current packet is pushed back
into the queue and the loop exits. The packet will be sent in the next slot of the node, one round
later. This back-pressure mechanism ensures that the transmitter thread stops \emph{offering}
data to the channel before its slot expires, regardless of queue depth or packet size.

It is worth being precise about what this mechanism does and does not control, because the
data plane runs over TCP. The transmitter releases each packet with a \texttt{send()} call,
which copies the data into the socket and transmit buffers of the kernel and returns; the kernel
then schedules the actual transmission on its own terms (TCP flow and congestion control,
then CSMA/CA medium access). The byte budget therefore bounds the \emph{volume} of data
handed to the kernel within a slot, which is the only quantity under the control of the framework;
it does not, and cannot, fix the exact instant at which the kernel puts each byte on the air.
A batch offered near the slot end may thus be transmitted by the kernel slightly into the
following slot. This is a deliberate consequence of operating the data plane over TCP, and the
byte budget is the lever that keeps the offered load slot-bounded; enforcing slot boundaries at
the wire level, below the socket, is discussed as future work in
Chapter~\ref{chap:conclusions}.


%% -----------------------------------------------------------------------
\section{Link Quality Monitoring}
\label{sec:pdr}

The routing decision in the work of Morais~\cite{morais2024synchronization} was based
solely on the binary connectivity matrix, treating all present links as equivalent regardless
of their reliability. This work introduces per-link Packet Delivery Ratio (PDR) monitoring
through the PDR module.

Each received \texttt{MSG\_DATA} packet carries a monotonically increasing sequence number
in the TDMA header. The PDR module counts received packets and compares the sequence
against the expected count, updating the estimate for the corresponding peer after every ten
received packets:

\begin{equation}
\mathrm{PDR}_{ij} = \frac{\text{packets received from peer } j}{%
  \text{packets expected from peer } j} \times 100
\label{eq:pdr-impl}
\end{equation}

The PDR value, an integer in $[0,\,100]$, is used as the edge weight in the MST computation.
Because the MST minimises total edge weight, a link with lower PDR (higher loss rate) has
a higher weight and is less likely to be selected when a more reliable alternative exists.
Figure~\ref{fig:quality} illustrates the three-step process: the full network topology with
PDR values is fed into Prim's MST algorithm, which produces a spanning tree retaining only
the highest-quality links; a DFS over that tree then derives the routing table with the
correct next-hop for each destination.

\begin{figure}[H]
\centering
\includegraphics[width=0.95\linewidth]{quality}
\caption{PDR-weighted MST routing: from network topology to routing table.}
\label{fig:quality}
\end{figure}



%% -----------------------------------------------------------------------
\section{End-to-End Packet Flow}
\label{sec:pkt-flow}

This section traces a video packet from source application to destination, covering both
the direct case (two nodes) and the relay case (three nodes).

Figure~\ref{fig:endtoend2} shows the direct case where source and destination are
neighbours. The source application writes an IP packet to \texttt{tun0}; the TUN reader
captures it and pushes it to the transmission queue; during its TDMA slot the transmitter
sends it over the TCP connection directly to the destination. The \texttt{tcp\_rx\_peer} thread of the destination receives it, calls \texttt{tun\_write()}, and the kernel
delivers it to the destination application via \texttt{tun0}.

\begin{figure}[H]
\centering
\includegraphics[width=1.0\textwidth]{EndtoEnd2}
\caption{Direct packet delivery between two neighbouring nodes.}
\label{fig:endtoend2}
\end{figure}

Figure~\ref{fig:endtoend} shows the three-node relay case, traced node by node below.

\begin{figure}[H]
\centering
\includegraphics[width=1.0\textwidth]{Endtoend}
\caption{End-to-end packet flow through a three-node relay topology.}
\label{fig:endtoend}
\end{figure}

At the source node, the video application writes IP packets addressed to the destination into \texttt{tun0}. The
TUN reader thread captures each packet, extracts the destination identifier from the IP
header, and pushes the packet into the transmission queue. During the TDMA slot of the source node, the transmitter thread pops the packet, calls \texttt{routing\_manager\_lookup()} to
obtain the next-hop peer ID, and sends the framed packet over the TCP connection to that
peer. If the slot boundary is reached, remaining packets are held for the following round.

At the relay node, the \texttt{tcp\_rx\_peer\_loop} thread receives the framed packet, updates the PDR estimate
for the source link, and inspects the destination IP address. Because the destination is not
the local node, the IP payload is written to \texttt{tun0} via \texttt{tun\_write()}. From
there the IP forwarding engine of the kernel takes over: the policy rule for traffic
ingressing on \texttt{tun0} directs the packet to table~200, whose \texttt{/32} host route
sends it straight out of \texttt{wlan0} to the physical \texttt{172.20.10.x} address of the
next hop. The relay therefore forwards the datagram at the network layer, without
re-queueing it for the TCP data plane of the framework, so the forwarded packet travels on the
air as a plain IP datagram rather than re-encapsulated framework traffic.

At the destination node, the forwarded IP datagram is received on \texttt{wlan0} and the kernel delivers
it to the application through \texttt{tun0}. The application thus receives the video stream
transparently, regardless of whether it arrived directly or through the relay.

\subsection{Generalisation to More Hops}
\label{subsec:more-hops}

Although the packet flow above was traced over a single relay, the forwarding
mechanism is independent of the number of hops. Every intermediate node performs
exactly the same operation. It receives the framed packet on its data plane,
writes the unmodified IP payload into \texttt{tun\textit{id}}, and lets the
kernel forward it out of \texttt{wlan0} toward the next physical hop through the
host route held in table~200. Nothing in this step depends on the position of the
node along the path or on the total path length, so the same mechanism carries a
datagram across \mbox{N1\,$\to$\,N2\,$\to$\,N3} and across an arbitrarily longer
chain \mbox{N1\,$\to$\,N2\,$\to$\,\ldots\,$\to$\,Nk} without any change.

The routing layer was designed for this from the start. Next hops are derived
from a Depth-First Search over the Minimum Spanning Tree of the full
connectivity matrix, which already spans every active node, so a longer path is
simply a tree with more vertices and produces additional routing-table entries
rather than a special case to be handled separately. No part of the algorithm,
the data structures, or the kernel configuration assumes a three-node topology.

What does grow with the hop count is governed by the TDMA structure rather than
by the routing design. Each additional relay adds one slot of store-and-forward
latency, because a packet that arrives during the slot of one node can only be
retransmitted in the next slot of that node, so the end-to-end delay increases by
roughly one frame period per hop. The single shared channel is also divided
among all active nodes, so the throughput available to any one flow falls as the
network grows, an effect inherent to a single-channel multi-hop TDMA chain
rather than to this implementation. Convergence after a link failure, by
contrast, remains bounded by the same node-aging timeout regardless of path
length, because each node detects a dead neighbour independently from its own
beacon reception. A systematic study of these effects on a larger testbed is
left as future work (Chapter~\ref{chap:conclusions}).



%% -----------------------------------------------------------------------
\section{Summary}
\label{sec:impl-summary}

This chapter presented the design and implementation of a Layer~3 routing system built on
top of the RA-TDMAs+ framework of Almeida~\cite{almeida2023synchronization}. The work
introduces four main contributions relative to the previous approach of
Morais~\cite{morais2024synchronization}.

First, the communication stack is separated into a UDP control plane for topology beacons
and a TCP data plane for application traffic. This separation ensures reliable, in-order
delivery of video frames without affecting the beacon mechanism, and is the primary
architectural motivation for operating at Layer~3.

Second, forwarding is performed at the Network layer through the IP
forwarding engine of the Linux kernel, making the relay path visible to standard diagnostic tools and
compatible with the networking subsystems of the operating system. A dual routing table
architecture separates application-originated traffic from relay-forwarded traffic, ensuring
correct forwarding without routing loops.

Third, per-link Packet Delivery Ratio monitoring is introduced as an edge weight in the
MST computation, so that the routing algorithm prefers reliable links over weak ones when
alternatives exist. This was not considered in the previous work, where all nominally
connected links were treated as equivalent.

Fourth, a per-slot byte budget with back-pressure ensures that the transmitter never
overruns its TDMA slot allocation, protecting the synchronisation quality of the entire
network.

Together these contributions produce a routing system that is observable, reliable, and
aware of link quality, while reusing the MST-based proactive algorithm and the
Topology-Tracking infrastructure already established by the framework. The following
chapter evaluates this system experimentally on physical robotic hardware.
